\section{Đạo đức nghề nghiệp và bảo đảm tính toàn vẹn của dữ liệu chất lượng}

\subsection{Chuẩn mực đạo đức nghề nghiệp trong quản lý chất lượng giáo dục}

Công tác quản lý chất lượng và kiểm định chất lượng giáo dục nghề nghiệp không đơn thuần là một chuỗi các thao tác kỹ thuật hay thủ tục hành chính, mà là hoạt động chuyên môn mang tính chuẩn mực đạo đức sâu sắc, gắn liền với niềm tin của xã hội và người sử dụng lao động đối với sản phẩm giáo dục. Bốn nguyên tắc đạo đức nghề nghiệp cốt lõi chi phối toàn bộ quá trình thu thập, xử lý dữ liệu minh chứng và công bố kết quả kiểm định.

Nguyên tắc thứ nhất là tính khách quan. Mọi hoạt động thu thập thông tin, đánh giá, chấm điểm và đưa ra kết luận kiểm định đều phải dựa trên các bằng chứng thực tế có thể kiểm chứng được một cách độc lập. Người làm công tác đảm bảo chất lượng phải giữ vững tính độc lập chuyên môn, tuyệt đối không để các áp lực về chỉ tiêu thành tích, định kiến cá nhân hay lợi ích cục bộ của đơn vị chi phối kết quả đánh giá.

Nguyên tắc thứ hai là tính trung thực. Dữ liệu chất lượng được xem là tấm gương phản chiếu chân thực mọi khía cạnh trong hoạt động vận hành của cơ sở giáo dục nghề nghiệp. Nhà trường nghiêm cấm tuyệt đối mọi hành vi ngụy tạo hồ sơ, hợp thức hóa minh chứng hồi cứu, can thiệp chỉnh sửa dữ liệu khảo sát gốc hoặc báo cáo sai lệch kết quả điều tra để tạo ra bức tranh chất lượng hoàn hảo giả tạo.

Nguyên tắc thứ ba là tính liêm chính học thuật. Nguyên tắc này đòi hỏi sự tôn trọng tuyệt đối tính chính xác của dữ liệu nghiên cứu và điều tra xã hội học. Toàn bộ thông tin thu thập từ người học, nhà giáo và doanh nghiệp phải được bảo lưu nguyên vẹn, không được phép cắt xén, chọn lọc thiên vị hay loại bỏ các phản hồi tiêu cực nhằm mục đích tô hồng thành tích.

Nguyên tắc thứ tư là trách nhiệm giải trình xã hội. Cơ sở giáo dục nghề nghiệp và đội ngũ làm công tác chất lượng mang trách nhiệm đạo đức lớn lao trước người học, gia đình và thị trường lao động. Việc công bố sai lệch tỷ lệ việc làm hay năng lực thực tế của người học sau tốt nghiệp là hành vi lừa dối người học và xã hội, làm sai lệch cán cân cung -- cầu nhân lực và làm suy giảm nghiêm trọng uy tín quốc gia của hệ thống giáo dục nghề nghiệp.

\subsection{Tình huống nghiên cứu điển hình: Áp lực chỉ tiêu và tính xác thực của dữ liệu}

\noindent\textit{(Ghi chú: Dưới đây là tình huống nghiên cứu điển hình giả định -- Case Study được xây dựng nhằm phục vụ phân tích, vận dụng và kiểm tra Chuẩn đầu ra 4 của học phần Quản lý chất lượng giáo dục nghề nghiệp. Các nhân vật, số liệu và tình tiết cụ thể mang tính chất minh họa học thuật, không phản ánh vụ việc có thật của bất kỳ cá nhân nào tại cơ sở đào tạo).}\par\vspace{0.2cm}

\subsubsection{Bối cảnh phát sinh vấn đề}

Giả định bối cảnh nghiên cứu: Vào thời điểm nhà trường chuẩn bị đón Đoàn đánh giá ngoài của Trung tâm Kiểm định chất lượng giáo dục nghề nghiệp độc lập đến khảo sát chính thức, phục vụ công nhận đạt tiêu chuẩn kiểm định cơ sở giáo dục nghề nghiệp theo quy định của Thông tư số 14/2024/TT-BLĐTBXH. Để bảo vệ uy tín thương hiệu của một cơ sở giáo dục nghề nghiệp công lập tiên phong tự chủ toàn diện, Hội đồng trường và Hiệu trưởng giao chỉ tiêu tại Mục tiêu chất lượng: tỷ lệ học sinh, sinh viên tốt nghiệp có việc làm phù hợp với ngành, nghề đào tạo trong thời gian mười hai tháng kể từ khi tốt nghiệp phải đạt từ 80\% trở lên.

Trong tình huống giả định này, chuyên viên phụ trách công tác tổng hợp dữ liệu kiểm định thuộc Phòng Đảm bảo chất lượng (tạm gọi là chuyên viên A) tiến hành rà soát, đối soát và làm sạch dữ liệu minh chứng phục vụ Tiêu chuẩn 8.2 về khảo sát việc làm sinh viên tốt nghiệp đối với khóa đào tạo gần nhất của Khoa Kỹ thuật Cơ khí -- Tự động hóa. Qua quá trình xử lý dữ liệu độc lập từ 350 cựu sinh viên của khoa, chuyên viên A phát hiện một bài toán thực tế đáng quan ngại: do ảnh hưởng của làn sóng tái cấu trúc sản xuất tại các doanh nghiệp vệ tinh trong khu vực, nhu cầu tuyển dụng kỹ thuật viên cơ khí sụt giảm đáng kể. Kết quả khảo sát thực tế cho thấy chỉ có 61,5\% sinh viên tìm được việc làm đúng hoặc phù hợp với chuyên môn đào tạo; 18,5\% sinh viên đang làm các công việc phổ thông tạm thời như chạy xe ôm công nghệ hoặc giao hàng trực tuyến; và 20,0\% sinh viên chưa có việc làm hoặc không duy trì liên lạc.

\subsubsection{Xung đột đạo đức và các áp lực thực tế}

Nhận thấy tỷ lệ việc làm thực tế 61,5\% của Khoa Kỹ thuật Cơ khí -- Tự động hóa sẽ kéo chỉ số bình quân toàn trường xuống dưới 75\%, làm không đạt cam kết mục tiêu chất lượng 80\% và khiến Tiêu chuẩn 8.2 có nguy cơ bị Đoàn đánh giá ngoài đánh giá Không đạt, Trưởng khoa Cơ khí đã chủ động đề xuất hai giải pháp can thiệp số liệu. Thứ nhất, gộp toàn bộ 18,5\% sinh viên làm công việc xe ôm công nghệ và giao hàng vào nhóm việc làm kỹ thuật viên dịch vụ vận tải và logistics công nghiệp nhằm hợp thức hóa thành việc làm phù hợp chuyên môn, qua đó nâng tỷ lệ có việc làm của khoa lên đúng 80,0\%. Thứ hai, loại bỏ bảy mươi phiếu khảo sát của những cựu sinh viên có phản hồi tiêu cực về việc khó tìm việc làm do trang thiết bị xưởng thực hành của trường thiếu các máy gia công điều khiển số hiện đại.

Trưởng khoa Cơ khí lập luận rằng sự sụt giảm việc làm là do bối cảnh suy thoái kinh tế khách quan, không xuất phát từ lỗi giảng dạy của khoa. Nếu báo cáo con số 61,5\%, uy tín tuyển sinh của trường sẽ bị ảnh hưởng nặng nề, trực tiếp đe dọa đến thu nhập và công việc của các giảng viên trong khoa. Thậm chí, trong cuộc họp giao ban nội bộ, một đồng chí Phó Hiệu trưởng phụ trách đào tạo cũng định hướng rằng bộ phận chất lượng cần có tư duy đồng hành bảo vệ thương hiệu trường, bảo đảm hồ sơ gửi ra ngoài phải đạt chuẩn để không làm tổn hại công sức chung của tập thể.

\begin{figure}[htbp]
\centering
\small
\setstretch{1.1}
\begin{tikzpicture}[
    root/.style={rectangle, draw=red!70, fill=red!10, thick, text width=13.6cm, align=center, rounded corners=3pt, minimum height=0.9cm, inner sep=6pt},
    choice1/.style={rectangle, draw=orange!70, fill=orange!5, thick, text width=6.3cm, align=left, rounded corners=3pt, inner sep=6pt},
    choice2/.style={rectangle, draw=teal!70, fill=teal!5, thick, text width=6.3cm, align=left, rounded corners=3pt, inner sep=6pt},
    arrow/.style={->, >=stealth, thick, color=gray!80}
]
\node[root, anchor=north] (ROOT) at (0, 0) {\textbf{TÌNH THẾ XUNG ĐỘT ĐẠO ĐỨC NGHỀ NGHIỆP CỦA CHUYÊN VIÊN CHẤT LƯỢNG}};
\node[choice1, anchor=north] (C1) at (-3.6, -1.8) {\textbf{LỰA CHỌN 1: THỎA HIỆP}\\[0.1cm] $\bullet$ Gộp việc phổ thông thành Logistics để đạt 80\%.\\[0.05cm] $\bullet$ Hủy các phiếu phản hồi tiêu cực.\\[0.1cm] \textbf{Hệ quả:} Báo cáo đạt chuẩn trước mắt nhưng vi phạm liêm chính và che giấu sự thật.};
\node[choice2, anchor=north] (C2) at (3.6, -1.8) {\textbf{LỰA CHỌN 2: BẢO VỆ DỮ LIỆU}\\[0.1cm] $\bullet$ Giữ nguyên tỷ lệ 61,5\% việc làm đúng chuyên môn.\\[0.05cm] $\bullet$ Bảo lưu 100\% dữ liệu gốc.\\[0.1cm] \textbf{Hệ quả:} Chịu áp lực nội bộ nhưng phản ánh đúng thực trạng để cải tiến thực chất.};

\coordinate (FORK) at (0, -1.35);
\draw[thick, color=gray!80] (ROOT.south) -- (FORK);
\draw[arrow] (FORK) -| (C1.north);
\draw[arrow] (FORK) -| (C2.north);
\end{tikzpicture}
\caption{\textbf{Sơ đồ phân tích tình thế xung đột đạo đức nghề nghiệp}}
\end{figure}

\subsection{Đánh giá đa chiều về rủi ro pháp lý, đạo đức và trách nhiệm xã hội}

Để đưa ra cách xử lý thấu đáo, người làm công tác quản lý chất lượng cần phân tích tình huống trên ba phương diện nền tảng.

Xét về phương diện đạo đức nghề nghiệp và liêm chính học thuật, hành vi tái phân loại việc làm xe ôm công nghệ thành kỹ thuật viên vận tải logistics và loại bỏ phiếu khảo sát tiêu cực là hành vi ngụy tạo dữ liệu và gian lận học thuật nghiêm trọng. Nếu chuyên viên đảm bảo chất lượng đồng thuận với phương án này, bộ phận chất lượng sẽ đánh mất hoàn toàn vai trò người bảo vệ chuẩn mực khách quan, biến hệ thống đảm bảo chất lượng thành công cụ hợp thức hóa thành tích ảo.

Xét về phương diện pháp lý và rủi ro kiểm định, quy trình đánh giá ngoài theo Thông tư số 14/2024/TT-BLĐTBXH quy định chặt chẽ việc kiểm định viên tiến hành phỏng vấn ngẫu nhiên cựu người học và xác minh chéo với người sử dụng lao động. Khi hành vi thao túng minh chứng bị phát hiện, Trường Cao đẳng Kỹ Nghệ II sẽ phải đối mặt với các chế tài xử phạt hành chính nghiêm khắc, bị hủy bỏ kết quả kiểm định và bị công khai sai phạm trên cổng thông tin của Bộ Lao động -- Thương binh và Xã hội. Tổn thất về mặt thương hiệu và uy tín tích lũy gần năm mươi năm của nhà trường là không thể cứu vãn.

Xét về phương diện trách nhiệm xã hội đối với người học, việc che đậy con số 61,5\% sẽ khiến nhà trường đánh mất cơ hội nhận diện khoảng trống công nghệ trong chương trình đào tạo ngành Cơ khí. Hậu quả lâu dài là các thế hệ người học tiếp theo sẽ tiếp tục tốt nghiệp với những kỹ năng chưa bắt kịp thực tiễn sản xuất công nghiệp hiện đại, đẩy gánh nặng việc làm lên vai người học và xã hội.

\subsection{Giải pháp và lộ trình xử lý xung đột của nhà quản lý chất lượng}

Trước áp lực thỏa hiệp, chuyên viên phụ trách chất lượng cùng Lãnh đạo Phòng Đảm bảo chất lượng kiên định thực hiện Quy trình năm bước xử lý xung đột chuẩn mực, kết hợp hài hòa giữa sự liêm chính về nguyên tắc với tinh thần xây dựng và giải pháp đồng hành.

\begin{figure}[htbp]
\centering
\small
\setstretch{1.1}
\begin{tikzpicture}[
    stepbox/.style={rectangle, draw=blue!70, fill=blue!5, thick, text width=12.5cm, align=left, rounded corners=3pt, minimum height=0.9cm},
    arrow/.style={->, >=stealth, thick, color=blue!80}
]
\node[stepbox] (S1) {\textbf{BƯỚC 1: ĐÓNG BĂNG VÀ BẢO LƯU DỮ LIỆU GỐC}\\ Sao lưu toàn bộ cơ sở dữ liệu khảo sát 350 sinh viên; xuất biên bản số hóa có chữ ký số.};
\node[stepbox, below=0.35cm of S1] (S2) {\textbf{BƯỚC 2: PHÂN TÍCH RỦI RO PHÁP LÝ VÀ CHUẨN MỰC KIỂM ĐỊNH}\\ Đối chiếu quy định Thông tư 14/2024; lập văn bản phân tích hệ lụy nếu làm sai lệch minh chứng.};
\node[stepbox, below=0.35cm of S2] (S3) {\textbf{BƯỚC 3: ĐỐI THOẠI CHUYÊN MÔN VỚI LÃNH ĐẠO KHOA CƠ KHÍ}\\ Thuyết phục khoa nhìn nhận số liệu 61,5\% là căn cứ khoa học để đề xuất mua sắm máy móc mới.};
\node[stepbox, below=0.35cm of S3] (S4) {\textbf{BƯỚC 4: BÁO CÁO HIỆU TRƯỞNG VÀ ĐỀ XUẤT GIẢI PHÁP HỌC THUẬT}\\ Trình báo cáo trung thực kèm Đề án phân tích nguyên nhân và Kế hoạch cải tiến chất lượng.};
\node[stepbox, below=0.35cm of S4] (S5) {\textbf{BƯỚC 5: BẢO VỆ TÍNH TRUNG THỰC TRƯỚC ĐOÀN ĐÁNH GIÁ NGOÀI}\\ Công khai trung thực tỷ lệ 61,5\%; trình bày Kế hoạch khắc phục bài bản để tạo niềm tin chuyên môn.};

\draw[arrow] (S1) -- (S2);
\draw[arrow] (S2) -- (S3);
\draw[arrow] (S3) -- (S4);
\draw[arrow] (S4) -- (S5);
\end{tikzpicture}
\caption{\textbf{Quy trình năm bước xử lý xung đột đạo đức nghề nghiệp}}
\end{figure}

Tại Bước 1, chuyên viên chất lượng tiến hành đóng băng và bảo lưu toàn vẹn tập dữ liệu khảo sát gốc của 350 cựu sinh viên Cơ khí, bao gồm nhật ký thời gian gửi phiếu, địa chỉ mạng và các bản ghi phản hồi chi tiết. Dữ liệu được mã hóa và lưu trữ tại kho lưu trữ số của Phòng Đảm bảo chất lượng có chữ ký số xác thực, bảo đảm không một cá nhân nào có thể tùy tiện can thiệp hoặc thay đổi dữ liệu.

Tại Bước 2, Phòng Đảm bảo chất lượng soạn thảo Bản đánh giá rủi ro pháp lý và quy chuẩn kiểm định, viện dẫn rõ các quy định xử phạt vi phạm hành chính trong giáo dục nghề nghiệp và trách nhiệm giải trình minh chứng. Văn bản làm rõ rằng việc bị phát hiện gian lận dữ liệu sẽ gây ra thảm họa về uy tín lớn hơn nhiều so với việc nhận điểm chưa đạt ở một tiêu chuẩn cụ thể.

Tại Bước 3, Lãnh đạo Phòng Đảm bảo chất lượng chủ trì buổi đối thoại chuyên môn với Ban Chủ nhiệm Khoa Kỹ thuật Cơ khí -- Tự động hóa. Buổi làm việc tập trung làm rõ: tỷ lệ 61,5\% là dữ liệu thực tế phản ánh bối cảnh suy thoái kinh tế và khoảng cách về thiết bị công nghệ của trường, không phải là cơ sở để quy kết trách nhiệm cá nhân thầy cô. Thay vì tìm cách che giấu, khoa cần sử dụng chính dữ liệu xác thực này làm luận cứ thuyết phục để kiến nghị Hiệu trưởng phê duyệt gói đầu tư trang bị trung tâm gia công điều khiển số (CNC) hiện đại và tổ chức các lớp bồi dưỡng kỹ năng số ngắn hạn cho cựu sinh viên chưa có việc làm.

Tại Bước 4, Phòng Đảm bảo chất lượng gửi Báo cáo giải trình chính thức lên Hiệu trưởng, kiên định quan điểm giữ nguyên số liệu trung thực 61,5\% trong Báo cáo tự đánh giá chính thức của trường. Đồng thời, báo cáo đính kèm một Kế hoạch hành động cải tiến chất lượng đột phá, bao gồm cam kết ký kết hợp tác tuyển dụng bổ sung với các hiệp hội doanh nghiệp cơ khí tại Thành phố Hồ Chí Minh và mở các khóa đào tạo bổ sung năng lực số nhằm nâng tỷ lệ việc làm lên 85\% trong sáu tháng tiếp theo.

Tại Bước 5, khi làm việc với Đoàn đánh giá ngoài, nhà trường công khai minh bạch số liệu 61,5\% tại Tiêu chuẩn 8.2 và chủ động giải trình sâu sắc về bối cảnh cùng Kế hoạch hành động khắc phục bài bản. Thực tiễn kiểm định chứng minh rằng các đoàn chuyên gia luôn đánh giá cao những cơ sở giáo dục nghề nghiệp dũng cảm nhìn thẳng vào sự thật, có năng lực nhận diện điểm nghẽn và đưa ra lộ trình cải tiến thuyết phục, hơn là một bộ hồ sơ hoàn hảo nhưng thiếu tính xác thực.

\subsection{Cơ chế kiểm soát nội bộ và bảo đảm tính trung thực của dữ liệu}

Để ngăn chặn tận gốc các áp lực thỏa hiệp và bảo vệ tính trung thực của hệ thống dữ liệu chất lượng lâu dài, Trường Cao đẳng Kỹ Nghệ II thiết lập hai cơ chế liêm chính mang tính thể chế.

Cơ chế thứ nhất là Quy chế bảo vệ người phản ánh sai phạm. Nhà trường thiết lập kênh tiếp nhận thông tin trực tuyến bảo mật độc lập, bao gồm hòm thư điện tử mã hóa và đường dây nóng do Trưởng Ban Thanh tra nhân dân và Đại diện Lãnh đạo về chất lượng trực tiếp quản lý. Nhà trường cam kết bảo vệ tuyệt đối danh tính, vị trí việc làm và quyền lợi chính đáng của mọi cán bộ, giảng viên, nhân viên và người học dũng cảm thông báo các hành vi thao túng hồ sơ, làm đẹp số liệu kiểm định hoặc trù dập người học. Mọi hành vi đe dọa, cản trở hoặc trù dập người phản ánh đều bị xử lý nghiêm minh theo quy định kỷ luật cao nhất.

Cơ chế thứ hai là Quy trình kiểm toán dữ liệu chất lượng độc lập. Trước khi Báo cáo tự đánh giá của nhà trường được ký ban hành chính thức, toàn bộ cơ sở dữ liệu minh chứng của năm mươi tiêu chuẩn bắt buộc phải trải qua một đợt kiểm toán dữ liệu độc lập do Tổ Kiểm toán chất lượng (gồm các chuyên gia kiểm định độc lập ngoài trường và đại diện doanh nghiệp đối tác) thực hiện. Tổ kiểm toán tiến hành đối soát ngẫu nhiên tối thiểu 10\% mẫu dữ liệu khảo sát việc làm và sự hài lòng bằng phương thức gọi điện xác minh trực tiếp. Nếu phát hiện tỷ lệ sai lệch vượt quá 3\%, toàn bộ hồ sơ của đơn vị đó sẽ bị yêu cầu làm sạch lại và người đứng đầu đơn vị phải chịu trách nhiệm giải trình trước Hội đồng trường.
